\section{Introducción}

Con el backend completamente implementado y una API REST funcional disponible para 
todos los módulos del sistema, el presente hito aborda el desarrollo del frontend. 
A diferencia del hito anterior, cuyo foco era la lógica de negocio y la persistencia 
de datos, este hito se centra en la experiencia del usuario: cómo se presenta la 
información, cómo se navega entre vistas y cómo se garantiza que cada rol acceda 
únicamente a lo que le corresponde.

El desarrollo se organiza en seis fases progresivas, donde cada una deja las 
condiciones preparadas para que la siguiente pueda avanzar sin fricciones. Esta 
estrategia permite mantener en todo momento una aplicación funcional y coherente, 
evitando estados intermedios rotos o incompletos. Las primeras fases establecen 
la base de autenticación, navegación y control de acceso; las intermedias 
construyen las vistas de gestión de inventario; y las últimas implementan el flujo 
de préstamos y refinan la experiencia visual y de usabilidad del sistema completo.

Un aspecto relevante que este hito pone de manifiesto es la naturaleza iterativa 
del desarrollo full-stack: a medida que el frontend avanzaba, surgieron necesidades 
que requirieron ajustes en el backend, como la corrección de relaciones 
\texttt{ON DELETE CASCADE} en la base de datos. Estos ajustes, gestionados mediante 
migraciones con Alembic, evidencian que ambas capas del sistema no se desarrollan 
de forma completamente independiente, sino que se retroalimentan a lo largo del 
proceso.